%-------------------------------------------------------------------------------
% Resume — Quantitative research roles (Citadel / Two Sigma / Jane Street, etc.)
% Zeyu (Steven) Zhang.  Compile with: tectonic resume-quant.tex
%-------------------------------------------------------------------------------

\documentclass[10.5pt]{article}

\usepackage[letterpaper, margin=1.5cm, bottom=1.6cm]{geometry}
\usepackage[T1]{fontenc}
\usepackage[utf8]{inputenc}
\usepackage{XCharter}
\usepackage[dvipsnames]{xcolor}
\usepackage{enumitem}
\usepackage{titlesec}
\usepackage{tabularx}
\usepackage{fontawesome5}
\usepackage{hyperref}

\definecolor{accent}{HTML}{1F3A6E}   % deep navy — conservative palette for finance
\definecolor{bodygray}{HTML}{2B2B2B}
\definecolor{metagray}{HTML}{5A5A5A}
\colorlet{chipbg}{accent!10}
\colorlet{rulegray}{black!25}

\hypersetup{colorlinks=true, urlcolor=accent, linkcolor=accent,
  pdftitle={Zeyu Zhang - Resume (Quantitative Research)}, pdfauthor={Zeyu Zhang}}

\color{bodygray}
\pagestyle{empty}
\setlength{\parindent}{0pt}
\raggedbottom

% Section title: small caps in accent, two-tone rule underneath
\titleformat{\section}
  {\large\scshape\color{accent}}
  {}{0em}{}
  [\vspace{-6pt}{\color{accent}\rule[0.6pt]{2.4cm}{1.6pt}}{\color{rulegray}\rule[1.15pt]{\dimexpr\textwidth-2.4cm\relax}{0.5pt}}]
\titlespacing*{\section}{0pt}{14pt}{8pt}

\newcommand{\chip}[1]{{\setlength{\fboxsep}{2.2pt}\colorbox{chipbg}{\scriptsize\bfseries\textcolor{accent}{#1}}}}

\newcommand{\impact}[1]{\par\vspace{3pt}{\small \textcolor{accent}{\faChartLine}\;\; \textbf{#1}}\par}

% Story lead: short narrative with accent bar (mirrors the website's research stories)
\newcommand{\story}[1]{%
  \par\vspace{3pt}
  \noindent{\color{accent}\vrule width 2pt}\hspace{8pt}%
  \begin{minipage}[t]{\dimexpr\textwidth-13pt\relax}
    {\small\itshape #1}
  \end{minipage}\par\vspace{3pt}}

% Bullet angle label
\newcommand{\facet}[1]{\textcolor{accent}{\textbf{#1}}}

\newcommand{\entry}[4]{%
  \begin{tabularx}{\textwidth}{@{}X@{}r@{}}
    \textbf{#1} & \textbf{#2} \\
    {\itshape\color{metagray}#3} & {\itshape\color{metagray}#4} \\
  \end{tabularx}\vspace{2pt}}

% Publication: {title}{chips}{authors + links}
\newcommand{\pub}[3]{%
  \item \textbf{#1}\\[1pt]
  {\footnotesize #3}\hfill #2}

\newenvironment{bullets}
  {\begin{itemize}[leftmargin=1.35em, topsep=2.5pt, itemsep=2pt, parsep=0pt,
                   label=\textcolor{accent}{\footnotesize$\bullet$}]\small}
  {\end{itemize}}

\newenvironment{publist}
  {\begin{itemize}[leftmargin=1.35em, topsep=2.5pt, itemsep=7pt, parsep=0pt,
                   label=\textcolor{accent}{\footnotesize$\bullet$}]\small}
  {\end{itemize}}

\newcommand{\me}{\textbf{Z.\ Zhang}}

%===============================================================================
\begin{document}

%--------------------------- Header --------------------------------------------
\begin{center}
  {\Huge\scshape Zeyu (Steven) Zhang}\\[6pt]
  {\color{metagray}\itshape Ph.D.\ Student in Statistics and Data Science, Northwestern University}\\[6pt]
  {\small
    \faEnvelope\ \href{mailto:zeyuzhang2028@u.northwestern.edu}{zeyuzhang2028@u.northwestern.edu}
    \;$|$\;
    \faGlobe\ \href{https://zeyuzhang1901.github.io}{zeyuzhang1901.github.io}
    \;$|$\;
    \faGithub\ \href{https://github.com/ZeyuZhang1901}{ZeyuZhang1901}
    \;$|$\;
    \faGraduationCap\ \href{https://scholar.google.com/citations?user=dr-xetEAAAAJ}{Google Scholar}
  }
\end{center}

%--------------------------- Summary --------------------------------------------
\section{Summary}
Statistics \& Data Science Ph.D.\ student at Northwestern (GPA 3.95/4.00) researching
\textbf{LLM backtesting}: evaluating the forecasting ability of large language models on
historical data, and detecting, attributing, and eliminating the \textbf{look-ahead bias}
that silently inflates backtest scores when a pretrained model has already ``read the future.''
First-author publication at NeurIPS, three first-author papers under review or in submission
(NeurIPS 2026, EMNLP 2026, TMLR), and an AISTATS 2026 Spotlight. Strong foundation in probability, statistics, optimization, and
offline reinforcement learning; fluent in Python.

%--------------------------- Education ------------------------------------------
\section{Education}
\entry{Northwestern University \; {\normalfont\small\itshape\color{metagray}Ph.D.\ in Statistics and Data Science, GPA 3.95/4.00}}
  {Evanston, IL}
  {Advisor: Prof.\ Bradly C.\ Stadie; expected graduation June 2028}{Sept 2023 -- Present}
\begin{bullets}
  \item Ph.D.\ coursework (all A): Statistical Theory \& Methods I--III, Applied Probability,
        Mathematics of Machine Learning, Reinforcement Learning, Mathematical Optimization,
        Analysis of Financial Data, Advanced Regression, Deep Learning.
\end{bullets}
\vspace{5pt}
\entry{University of Science and Technology of China \; {\normalfont\small\itshape\color{metagray}B.Eng.\ EE; Minor in AI, GPA 3.93/4.30}}
  {Hefei, China}
  {Rank 5/213 in School of Information Science and Technology}{Sept 2019 -- Jun 2023}
\begin{bullets}
  \item Coursework: Stochastic Processes (97/100), Linear Algebra (97/100), Information Theory (93/100),
        Probability Theory (90/100), Machine Learning (93/100).
\end{bullets}

%--------------------------- Research Experience ---------------------------------
\section{Research Experience}

\entry{LLM Backtesting \& Look-Ahead Bias: Theory, Detection, and Elimination}
  {Northwestern University}
  {Advisor: Prof.\ Bradly C.\ Stadie \quad Three first-author papers in submission (NeurIPS / EMNLP 2026, TMLR)}{Sept 2025 -- Present}
\story{Backtesting is how every forecasting system is validated---and LLM forecasters, increasingly
  deployed for financial prediction, are no exception. But a pretrained LLM has already read the future:
  post-cutoff knowledge from pretraining leaks into its reasoning, a model-internal form of
  \textbf{look-ahead bias} that no date filter on the data can remove, silently inflating backtest scores.
  My research treats this as one program: a \textbf{statistical theory} of why the inflation is
  systematic and how to subtract it from the score, an \textbf{interpretable audit} and inference-time
  defense, and a \textbf{training method} that instills point-in-time discipline.}
\begin{bullets}
  \item \facet{Formalize the mechanism.} Modeled leakage causally and proved the inflation is systematic,
        not random: it obeys a stakes $\times$ extraction law and concentrates where the outcome surprised
        the contemporaneous crowd and the model absorbed it. Partial leakage is disproportionately
        rewarded: absorbing half the leaked signal already yields 75\% of the full inflation.
  \item \facet{Recover true forecasting skill.} Proved the inflation is \emph{non-identifiable} from
        backtest scores alone---entangled with genuine skill and recency, so a flat pre/post performance
        profile does \emph{not} certify a clean backtest---and that one practical external reference
        restores identification (\textbf{difference-in-differences} against a clean control model, or
        \textbf{regression discontinuity} at the training cutoff), yielding a \textbf{leakage-adjusted
        score} with confidence intervals. Estimators recover leakage planted in twin models in dose,
        location, and per-question profile, return null on the clean twin, and, deployed on frontier
        models over a 1{,}646-question prediction-market panel (Polymarket, Metaculus, Manifold, INFER),
        detect one cutoff-localized signature while clearing five models whose apparent forecasting edge
        was recency alone.
  \item \facet{Audit any black box.} \textbf{Shapley-DCLR}: decomposes each rationale into atomic claims
        and attributes each claim's decision impact with \textbf{Shapley values} (near-linear
        leverage-score estimation), measuring what fraction of the reasoning that \emph{actually drives} a
        prediction depends on future information.
  \item \facet{Block at inference.} \textbf{TimeSPEC}, a point-in-time evaluation architecture
        (date-filtered retrieval + claim-level supervision + closed-world aggregation) that blocks leaked
        claims at the output level with no retraining; evaluated on \textbf{2{,}769 forecasting instances}
        across three LLMs.
  \item \facet{Train it out.} Temporal compliance is \emph{instance-specific}---the same fact is legitimate
        for one cutoff date and a violation for another---so unlearning cannot fix it. \textbf{TEMPO}
        instead trains the discipline: constrained RL post-training with a two-mode reward that eliminates
        leaked claims as a hard constraint before optimizing predictive performance; \textbf{proved}
        monotonic descent of leakage and linear convergence to the leak-free optimum (PL condition).
  \item \facet{Point-in-time data.} Constructed three \textbf{point-in-time benchmarks} ($\sim$3{,}500
        instances) with per-instance cutoff dates and zero-knowledge controls: ranking same-sector stocks
        by six-month total return (incl.\ the COVID regime shift), athlete contract values, and federal
        case outcomes.
\end{bullets}
\impact{Leakage: retrieval baselines 13.1\% $\rightarrow$ 1.8\%; trained models 2--13\% $\rightarrow$
  0.6--3.7\%, predictive performance +6--13\%, zero inference overhead.}

\vspace{7pt}
\entry{Counterfactual Learning from Logged Feedback: Off-Policy Learning to Rank}
  {Princeton University (remote)}
  {Mentors: Prof.\ Mengdi Wang, Prof.\ Huazheng Wang \quad Published at NeurIPS 2023}{May 2022 -- Dec 2023}
\story{Learning to rank from logged user feedback is a counterfactual problem: clicks are biased by how
  results were displayed, and existing corrections each assume a specific, known user-behavior model.}
\begin{bullets}
  \item Cast ranking under general stochastic user-behavior (click) models as a Markov Decision Process,
        enabling offline RL (CQL, SAC, BCQ) to learn optimal rankings from biased logged data
        \emph{without} explicit debiasing.
  \item Statistically significant state-of-the-art on Yahoo!\ LETOR and MSLR-WEB10K, robust across
        position-based, cascade, and DCM behavior models.
\end{bullets}

\vspace{7pt}
\entry{Sensitivity Analysis of Black-Box Models via Local Linear Surrogates (LAMP)}
  {Northwestern University}
  {Advisor: Prof.\ Bradly C.\ Stadie \quad AISTATS 2026 \textbf{Spotlight}}{Jan 2024 -- Jan 2025}
\story{Before trusting a black-box model's output, one should ask whether the factors it cites actually
  move its predictions---a question of local sensitivity, answerable by regression.}
\begin{bullets}
  \item Co-developed a perturbation-based method that regresses a black-box model's re-queried predictions
        on perturbed factor weights, recovering a local linear approximation of its decision surface---no
        gradients or internals required; validated against human and clinical-expert judgments.
\end{bullets}

%--------------------------- Publications ---------------------------------------
\section{Publications \& Preprints}
\begin{publist}
  \pub{Temporal Leakage in LLM Backtesting: Measurement, Validation, and Adjusted Scores}
    {\chip{In submission, TMLR}}
    {\me, B.\ C.\ Stadie \; \href{https://github.com/ZeyuZhang1901/Temporal-Leakage-Backtesting}{[Code]}}
  \pub{TEMPO: Temporal Enforcement via Mode-Separated Policy Optimization for Trustworthy LLM Backtesting}
    {\chip{Under review, NeurIPS 2026}}
    {\me, B.\ C.\ Stadie \; \href{https://arxiv.org/abs/2605.18843}{[arXiv:2605.18843]}}
  \pub{All Leaks Count, Some Count More: Interpretable Temporal Contamination Detection and Mitigation in LLM Backtesting}
    {\chip{Under review, EMNLP 2026}}
    {\me, R.\ Chen, B.\ C.\ Stadie \; \href{https://arxiv.org/abs/2602.17234}{[arXiv:2602.17234]}}
  \pub{LAMP: Extracting Locally Linear Decision Surfaces from LLM World Models}
    {\chip{AISTATS 2026}\,\chip{Spotlight}}
    {R.\ Chen, Y.\ Ko, \me, C.\ Cho, S.\ Chung, M.\ Giuffr\'e, D.\ L.\ Shung, B.\ C.\ Stadie \;
     \href{https://arxiv.org/abs/2505.11772}{[arXiv:2505.11772]}}
  \pub{Unified Off-Policy Learning to Rank: a Reinforcement Learning Perspective}
    {\chip{NeurIPS 2023}}
    {\me, Y.\ Su, H.\ Yuan, Y.\ Wu, R.\ Balasubramanian, Q.\ Wu, H.\ Wang, M.\ Wang \;
     \href{https://arxiv.org/abs/2306.07528}{[arXiv:2306.07528]}}
\end{publist}

%--------------------------- Honors & Skills -------------------------------------
\section{Honors, Skills \& Teaching}
\begin{bullets}
  \item \textbf{Honors:} China National Scholarship, Ministry of Education (top 1\%), 2020;
        National Undergraduate Electronic Design Contest, 2nd Prize nationally / 1st in Anhui Province, 2021.
  \item \textbf{Programming:} Python (NumPy, Pandas, PyTorch), R, C/C++, Bash, Git, \LaTeX.
  \item \textbf{Quantitative:} probability \& stochastic processes, statistical inference \& experiment design,
        causal identification (difference-in-differences, regression discontinuity, partial identification),
        Shapley-value attribution, offline reinforcement learning, ML evaluation under data contamination.
  \item \textbf{Teaching:} TA for six quarters at Northwestern Statistics (probability \& statistics,
        multivariate analysis, data science with Python), 2024--2026.
\end{bullets}

\end{document}
